% Section: Tiền xử lý văn bản (Text Preprocessing) — nhánh lexical (TF-IDF / tf-isf)
% Không cần gói ngoài (chỉ dùng itemize + toán cơ bản).
% Dùng: % Section: Tiền xử lý văn bản (Text Preprocessing) — nhánh lexical (TF-IDF / tf-isf)
% Không cần gói ngoài (chỉ dùng itemize + toán cơ bản).
% Dùng: % Section: Tiền xử lý văn bản (Text Preprocessing) — nhánh lexical (TF-IDF / tf-isf)
% Không cần gói ngoài (chỉ dùng itemize + toán cơ bản).
% Dùng: % Section: Tiền xử lý văn bản (Text Preprocessing) — nhánh lexical (TF-IDF / tf-isf)
% Không cần gói ngoài (chỉ dùng itemize + toán cơ bản).
% Dùng: \input{report_preprocessing_section}

\section{Tiền xử lý văn bản (Text Preprocessing)}

Tầng đối chiếu (alignment) làm việc trên đơn vị \emph{câu}: mỗi câu của tài liệu nghi vấn được
so khớp với các câu của tài liệu nguồn. Do đó, sau bước truy hồi, văn bản cần được tách thành
câu trước khi tính độ tương đồng tf-isf. Bước tiền xử lý này được cài đặt trong
\texttt{scripts/alignment/align\_core.py}, với ràng buộc thiết kế cốt lõi là \emph{bảo toàn chính
xác vị trí ký tự}.

\subsection{Ràng buộc: chấm điểm ở mức ký tự}

Thước đo PlagDet đánh giá ở mức ký tự: một đoạn đạo văn được biểu diễn bằng khoảng
$[\,\mathrm{offset},\ \mathrm{offset}+\mathrm{length}\,)$ trong tài liệu, và điểm số phụ thuộc vào phần
giao ký tự giữa đoạn dự đoán và nhãn gold. Vì vậy, một sai lệch nhỏ ở offset câu sẽ làm lệch điểm
mà không phát sinh lỗi tường minh. Ràng buộc này loại trừ các bộ tách câu ``làm sạch'' văn bản
(chuẩn hoá khoảng trắng, bỏ ký tự, cắt đoạn), vì mọi thao tác như vậy đều phá vỡ ánh xạ offset
giữa câu và văn bản gốc.

\subsection{Bộ tách câu lát kín (tiling)}

Để thoả ràng buộc trên, bộ tách câu được thiết kế sao cho các câu \emph{lát kín} (tiling) toàn bộ
văn bản. Thuật toán dựa trên luật: quét văn bản, ngắt câu tại chuỗi ký tự kết thúc câu
\texttt{.!?} (một chuỗi dấu liên tiếp được gộp thành một ranh giới duy nhất), sau đó \emph{nuốt}
toàn bộ khoảng trắng theo sau vào câu vừa cắt. Phần văn bản còn lại sau dấu kết câu cuối cùng
(nếu có) tạo thành câu cuối.

Kết quả là một dãy câu, mỗi câu chỉ lưu cặp offset \texttt{(start, end)}, thoả các bất biến sau:

\begin{itemize}
  \item Câu đầu bắt đầu tại offset $0$; câu cuối kết thúc tại độ dài văn bản.
  \item Hai câu kề nhau khít nhau: câu thứ $i$ kết thúc đúng tại vị trí câu thứ $i{+}1$ bắt đầu
  --- không có khoảng hở, không chồng lấn.
  \item Khoảng trắng (kể cả xuống dòng) luôn thuộc về đúng một câu, nên việc ghép nối tất cả các
  câu tái tạo lại văn bản gốc \emph{từng ký tự}.
\end{itemize}

Bất biến này được kiểm chứng tự động bằng hàm \texttt{verify\_tiling}, khẳng định (assertion) rằng
dãy câu không hở, không chồng, và chuỗi ghép lại trùng khớp văn bản gốc. Nhờ đó, span đạo văn do
aligner sinh ra --- vốn được tạo bằng cách gộp các câu liên tiếp --- luôn có offset ký tự chính
xác, đảm bảo PlagDet không bị lệch điểm do sai số tách câu.

\subsection{Về việc không loại reference marker}

Ở nhánh lexical này, văn bản \emph{không} bị khoanh vùng nội dung (giữ nguyên phần đầu tài liệu và
danh mục tài liệu tham khảo) và \emph{không} loại các dấu trích dẫn (reference marker) như
\texttt{[12]} hay \texttt{(3-5)}. Đây là lựa chọn có chủ đích: việc cắt bỏ hoặc bỏ qua các đoạn văn
bản sẽ tạo ra khoảng hở trong dãy offset và phá vỡ tính lát kín, mâu thuẫn với yêu cầu bảo toàn ký
tự. Trên thực tế, các dấu trích dẫn rất ngắn và hiếm khi đạt ngưỡng để hình thành một span đạo văn
(aligner tf-isf yêu cầu độ tương đồng đủ cao trên một cụm câu liền mạch và độ dài tối thiểu), nên
việc giữ lại chúng gần như không ảnh hưởng đến độ chính xác, trong khi loại bỏ có thể gây rủi ro
lệch offset. Bước lọc reference marker chỉ được áp dụng ở biến thể aligner ngữ nghĩa, nơi mỗi câu
được nhúng độc lập và không có ràng buộc lát kín.


\section{Tiền xử lý văn bản (Text Preprocessing)}

Tầng đối chiếu (alignment) làm việc trên đơn vị \emph{câu}: mỗi câu của tài liệu nghi vấn được
so khớp với các câu của tài liệu nguồn. Do đó, sau bước truy hồi, văn bản cần được tách thành
câu trước khi tính độ tương đồng tf-isf. Bước tiền xử lý này được cài đặt trong
\texttt{scripts/alignment/align\_core.py}, với ràng buộc thiết kế cốt lõi là \emph{bảo toàn chính
xác vị trí ký tự}.

\subsection{Ràng buộc: chấm điểm ở mức ký tự}

Thước đo PlagDet đánh giá ở mức ký tự: một đoạn đạo văn được biểu diễn bằng khoảng
$[\,\mathrm{offset},\ \mathrm{offset}+\mathrm{length}\,)$ trong tài liệu, và điểm số phụ thuộc vào phần
giao ký tự giữa đoạn dự đoán và nhãn gold. Vì vậy, một sai lệch nhỏ ở offset câu sẽ làm lệch điểm
mà không phát sinh lỗi tường minh. Ràng buộc này loại trừ các bộ tách câu ``làm sạch'' văn bản
(chuẩn hoá khoảng trắng, bỏ ký tự, cắt đoạn), vì mọi thao tác như vậy đều phá vỡ ánh xạ offset
giữa câu và văn bản gốc.

\subsection{Bộ tách câu lát kín (tiling)}

Để thoả ràng buộc trên, bộ tách câu được thiết kế sao cho các câu \emph{lát kín} (tiling) toàn bộ
văn bản. Thuật toán dựa trên luật: quét văn bản, ngắt câu tại chuỗi ký tự kết thúc câu
\texttt{.!?} (một chuỗi dấu liên tiếp được gộp thành một ranh giới duy nhất), sau đó \emph{nuốt}
toàn bộ khoảng trắng theo sau vào câu vừa cắt. Phần văn bản còn lại sau dấu kết câu cuối cùng
(nếu có) tạo thành câu cuối.

Kết quả là một dãy câu, mỗi câu chỉ lưu cặp offset \texttt{(start, end)}, thoả các bất biến sau:

\begin{itemize}
  \item Câu đầu bắt đầu tại offset $0$; câu cuối kết thúc tại độ dài văn bản.
  \item Hai câu kề nhau khít nhau: câu thứ $i$ kết thúc đúng tại vị trí câu thứ $i{+}1$ bắt đầu
  --- không có khoảng hở, không chồng lấn.
  \item Khoảng trắng (kể cả xuống dòng) luôn thuộc về đúng một câu, nên việc ghép nối tất cả các
  câu tái tạo lại văn bản gốc \emph{từng ký tự}.
\end{itemize}

Bất biến này được kiểm chứng tự động bằng hàm \texttt{verify\_tiling}, khẳng định (assertion) rằng
dãy câu không hở, không chồng, và chuỗi ghép lại trùng khớp văn bản gốc. Nhờ đó, span đạo văn do
aligner sinh ra --- vốn được tạo bằng cách gộp các câu liên tiếp --- luôn có offset ký tự chính
xác, đảm bảo PlagDet không bị lệch điểm do sai số tách câu.

\subsection{Về việc không loại reference marker}

Ở nhánh lexical này, văn bản \emph{không} bị khoanh vùng nội dung (giữ nguyên phần đầu tài liệu và
danh mục tài liệu tham khảo) và \emph{không} loại các dấu trích dẫn (reference marker) như
\texttt{[12]} hay \texttt{(3-5)}. Đây là lựa chọn có chủ đích: việc cắt bỏ hoặc bỏ qua các đoạn văn
bản sẽ tạo ra khoảng hở trong dãy offset và phá vỡ tính lát kín, mâu thuẫn với yêu cầu bảo toàn ký
tự. Trên thực tế, các dấu trích dẫn rất ngắn và hiếm khi đạt ngưỡng để hình thành một span đạo văn
(aligner tf-isf yêu cầu độ tương đồng đủ cao trên một cụm câu liền mạch và độ dài tối thiểu), nên
việc giữ lại chúng gần như không ảnh hưởng đến độ chính xác, trong khi loại bỏ có thể gây rủi ro
lệch offset. Bước lọc reference marker chỉ được áp dụng ở biến thể aligner ngữ nghĩa, nơi mỗi câu
được nhúng độc lập và không có ràng buộc lát kín.


\section{Tiền xử lý văn bản (Text Preprocessing)}

Tầng đối chiếu (alignment) làm việc trên đơn vị \emph{câu}: mỗi câu của tài liệu nghi vấn được
so khớp với các câu của tài liệu nguồn. Do đó, sau bước truy hồi, văn bản cần được tách thành
câu trước khi tính độ tương đồng tf-isf. Bước tiền xử lý này được cài đặt trong
\texttt{scripts/alignment/align\_core.py}, với ràng buộc thiết kế cốt lõi là \emph{bảo toàn chính
xác vị trí ký tự}.

\subsection{Ràng buộc: chấm điểm ở mức ký tự}

Thước đo PlagDet đánh giá ở mức ký tự: một đoạn đạo văn được biểu diễn bằng khoảng
$[\,\mathrm{offset},\ \mathrm{offset}+\mathrm{length}\,)$ trong tài liệu, và điểm số phụ thuộc vào phần
giao ký tự giữa đoạn dự đoán và nhãn gold. Vì vậy, một sai lệch nhỏ ở offset câu sẽ làm lệch điểm
mà không phát sinh lỗi tường minh. Ràng buộc này loại trừ các bộ tách câu ``làm sạch'' văn bản
(chuẩn hoá khoảng trắng, bỏ ký tự, cắt đoạn), vì mọi thao tác như vậy đều phá vỡ ánh xạ offset
giữa câu và văn bản gốc.

\subsection{Bộ tách câu lát kín (tiling)}

Để thoả ràng buộc trên, bộ tách câu được thiết kế sao cho các câu \emph{lát kín} (tiling) toàn bộ
văn bản. Thuật toán dựa trên luật: quét văn bản, ngắt câu tại chuỗi ký tự kết thúc câu
\texttt{.!?} (một chuỗi dấu liên tiếp được gộp thành một ranh giới duy nhất), sau đó \emph{nuốt}
toàn bộ khoảng trắng theo sau vào câu vừa cắt. Phần văn bản còn lại sau dấu kết câu cuối cùng
(nếu có) tạo thành câu cuối.

Kết quả là một dãy câu, mỗi câu chỉ lưu cặp offset \texttt{(start, end)}, thoả các bất biến sau:

\begin{itemize}
  \item Câu đầu bắt đầu tại offset $0$; câu cuối kết thúc tại độ dài văn bản.
  \item Hai câu kề nhau khít nhau: câu thứ $i$ kết thúc đúng tại vị trí câu thứ $i{+}1$ bắt đầu
  --- không có khoảng hở, không chồng lấn.
  \item Khoảng trắng (kể cả xuống dòng) luôn thuộc về đúng một câu, nên việc ghép nối tất cả các
  câu tái tạo lại văn bản gốc \emph{từng ký tự}.
\end{itemize}

Bất biến này được kiểm chứng tự động bằng hàm \texttt{verify\_tiling}, khẳng định (assertion) rằng
dãy câu không hở, không chồng, và chuỗi ghép lại trùng khớp văn bản gốc. Nhờ đó, span đạo văn do
aligner sinh ra --- vốn được tạo bằng cách gộp các câu liên tiếp --- luôn có offset ký tự chính
xác, đảm bảo PlagDet không bị lệch điểm do sai số tách câu.

\subsection{Về việc không loại reference marker}

Ở nhánh lexical này, văn bản \emph{không} bị khoanh vùng nội dung (giữ nguyên phần đầu tài liệu và
danh mục tài liệu tham khảo) và \emph{không} loại các dấu trích dẫn (reference marker) như
\texttt{[12]} hay \texttt{(3-5)}. Đây là lựa chọn có chủ đích: việc cắt bỏ hoặc bỏ qua các đoạn văn
bản sẽ tạo ra khoảng hở trong dãy offset và phá vỡ tính lát kín, mâu thuẫn với yêu cầu bảo toàn ký
tự. Trên thực tế, các dấu trích dẫn rất ngắn và hiếm khi đạt ngưỡng để hình thành một span đạo văn
(aligner tf-isf yêu cầu độ tương đồng đủ cao trên một cụm câu liền mạch và độ dài tối thiểu), nên
việc giữ lại chúng gần như không ảnh hưởng đến độ chính xác, trong khi loại bỏ có thể gây rủi ro
lệch offset. Bước lọc reference marker chỉ được áp dụng ở biến thể aligner ngữ nghĩa, nơi mỗi câu
được nhúng độc lập và không có ràng buộc lát kín.


\section{Tiền xử lý văn bản (Text Preprocessing)}

Tầng đối chiếu (alignment) làm việc trên đơn vị \emph{câu}: mỗi câu của tài liệu nghi vấn được
so khớp với các câu của tài liệu nguồn. Do đó, sau bước truy hồi, văn bản cần được tách thành
câu trước khi tính độ tương đồng tf-isf. Bước tiền xử lý này được cài đặt trong
\texttt{scripts/alignment/align\_core.py}, với ràng buộc thiết kế cốt lõi là \emph{bảo toàn chính
xác vị trí ký tự}.

\subsection{Ràng buộc: chấm điểm ở mức ký tự}

Thước đo PlagDet đánh giá ở mức ký tự: một đoạn đạo văn được biểu diễn bằng khoảng
$[\,\mathrm{offset},\ \mathrm{offset}+\mathrm{length}\,)$ trong tài liệu, và điểm số phụ thuộc vào phần
giao ký tự giữa đoạn dự đoán và nhãn gold. Vì vậy, một sai lệch nhỏ ở offset câu sẽ làm lệch điểm
mà không phát sinh lỗi tường minh. Ràng buộc này loại trừ các bộ tách câu ``làm sạch'' văn bản
(chuẩn hoá khoảng trắng, bỏ ký tự, cắt đoạn), vì mọi thao tác như vậy đều phá vỡ ánh xạ offset
giữa câu và văn bản gốc.

\subsection{Bộ tách câu lát kín (tiling)}

Để thoả ràng buộc trên, bộ tách câu được thiết kế sao cho các câu \emph{lát kín} (tiling) toàn bộ
văn bản. Thuật toán dựa trên luật: quét văn bản, ngắt câu tại chuỗi ký tự kết thúc câu
\texttt{.!?} (một chuỗi dấu liên tiếp được gộp thành một ranh giới duy nhất), sau đó \emph{nuốt}
toàn bộ khoảng trắng theo sau vào câu vừa cắt. Phần văn bản còn lại sau dấu kết câu cuối cùng
(nếu có) tạo thành câu cuối.

Kết quả là một dãy câu, mỗi câu chỉ lưu cặp offset \texttt{(start, end)}, thoả các bất biến sau:

\begin{itemize}
  \item Câu đầu bắt đầu tại offset $0$; câu cuối kết thúc tại độ dài văn bản.
  \item Hai câu kề nhau khít nhau: câu thứ $i$ kết thúc đúng tại vị trí câu thứ $i{+}1$ bắt đầu
  --- không có khoảng hở, không chồng lấn.
  \item Khoảng trắng (kể cả xuống dòng) luôn thuộc về đúng một câu, nên việc ghép nối tất cả các
  câu tái tạo lại văn bản gốc \emph{từng ký tự}.
\end{itemize}

Bất biến này được kiểm chứng tự động bằng hàm \texttt{verify\_tiling}, khẳng định (assertion) rằng
dãy câu không hở, không chồng, và chuỗi ghép lại trùng khớp văn bản gốc. Nhờ đó, span đạo văn do
aligner sinh ra --- vốn được tạo bằng cách gộp các câu liên tiếp --- luôn có offset ký tự chính
xác, đảm bảo PlagDet không bị lệch điểm do sai số tách câu.

\subsection{Về việc không loại reference marker}

Ở nhánh lexical này, văn bản \emph{không} bị khoanh vùng nội dung (giữ nguyên phần đầu tài liệu và
danh mục tài liệu tham khảo) và \emph{không} loại các dấu trích dẫn (reference marker) như
\texttt{[12]} hay \texttt{(3-5)}. Đây là lựa chọn có chủ đích: việc cắt bỏ hoặc bỏ qua các đoạn văn
bản sẽ tạo ra khoảng hở trong dãy offset và phá vỡ tính lát kín, mâu thuẫn với yêu cầu bảo toàn ký
tự. Trên thực tế, các dấu trích dẫn rất ngắn và hiếm khi đạt ngưỡng để hình thành một span đạo văn
(aligner tf-isf yêu cầu độ tương đồng đủ cao trên một cụm câu liền mạch và độ dài tối thiểu), nên
việc giữ lại chúng gần như không ảnh hưởng đến độ chính xác, trong khi loại bỏ có thể gây rủi ro
lệch offset. Bước lọc reference marker chỉ được áp dụng ở biến thể aligner ngữ nghĩa, nơi mỗi câu
được nhúng độc lập và không có ràng buộc lát kín.
